% GENERATED FILE. Edit canonical lessons, then run npm run generate:books.
% canonical-source-hash: fnv1a64:55afcf243f642e47
% canonical-lessons: TA-C408-valaivu, TA-C408-carivu, TA-C408-erram, TA-C408-kavil, TA-C408-culal

\chapter{Bends and Slopes}
\label{ch:ta-bends-and-slopes}
\hlchaptermodality{408}

\begin{chapteropening}
\textbf{By the end of this chapter:} \emph{I can talk about bends, slopes and climbs, and things overturning and spinning.}

Complete the last lesson of chapter 408: I can talk about bends, slopes and climbs, and things overturning and spinning.
\end{chapteropening}

\section[\texorpdfstring{\ta{வளைவு} (vaḷaivu)}{வளைவு (vaḷaivu)}]{\ta{வளைவு} (vaḷaivu) — a curve; a bend}

\label{lesson:TA-C408-valaivu}



\begin{quote}
Before the new one: say the Tamil for a route, then the Tamil for to pull over.
\end{quote}

\subsection*{You'll want to know: \ta{வளைவு}}
\textbf{\ta{வளைவு}} — \emph{vaḷaivu} — \textquotedblleft{}a curve; a bend\textquotedblright{}.

\begin{quote}
\ta{இந்த} \ta{மலைச்} \ta{சாலையில்} \ta{நிறைய} \ta{வளைவுகள்} \ta{இருக்கின்றன}.
\end{quote}

\begin{quote}
\emph{inta malaic cālaiyil niṟaiya vaḷaivukaḷ irukkiṉṟaṉa.}
\end{quote}

\begin{quote}
There are many bends on this hill road.
\end{quote}

Plural: \ta{வளைவுகள்}, \emph{vaḷaivukaḷ}.

From \ta{வளை}, to bend, + the noun ending -\ta{வு}.

\begin{cousinweb}
\textbf{\ta{வளைவு}} (\emph{vaḷaivu}) is built from \ta{வளை}, \textquotedblleft{}to bend, to become crooked\textquotedblright{} (DEDR 5314), with the noun ending \emph{-vu}.
\end{cousinweb}

\begin{grammarlens}[title={What you've built}]
One more word for getting around.
\end{grammarlens}

\subsection*{Your turn}
\begin{itemize}
\raggedright
  \item \emph{Say it:} \emph{vaḷaivu}
  \item \emph{Say it:} \emph{vaḷaivu}, once more
  \item \emph{Say it:} \emph{inta malaic cālaiyil niṟaiya vaḷaivukaḷ irukkiṉṟaṉa.} — There are many bends on this hill road.
  \item \emph{Recall:} say the Tamil for a route, then the Tamil for to pull over, then say \emph{vaḷaivu} again
\end{itemize}

\begin{tcolorbox}[breakable,colback=teal!4,colframe=teal!35!black,title={Before you move on}]
What does \ta{வளைவு} mean? (\textquotedblleft{}A curve; a bend\textquotedblright{}.) Say it once more.
\end{tcolorbox}
\section[\texorpdfstring{\ta{சரிவு} (carivu)}{சரிவு (carivu)}]{\ta{சரிவு} (carivu) — a slope; a decline}

\label{lesson:TA-C408-carivu}



\begin{quote}
Before the new one: say the Tamil for to pull over, then the Tamil for a curve.
\end{quote}

\subsection*{You'll want to know: \ta{சரிவு}}
\textbf{\ta{சரிவு}} — \emph{carivu} — \textquotedblleft{}a slope; a decline\textquotedblright{}.

\begin{quote}
\ta{மலைச்} \ta{சரிவில்} \ta{ஆடுகள்} \ta{மேய்ந்துகொண்டிருந்தன}.
\end{quote}

\begin{quote}
\emph{malaic carivil āṭukaḷ mēyntukoṇṭiruntaṉa.}
\end{quote}

\begin{quote}
Goats were grazing on the hillside.
\end{quote}

From \ta{சரி}, to slide down, + -\ta{வு}; \ta{விலைச்} \ta{சரிவு} is a fall in prices.

\begin{cousinweb}
\textbf{\ta{சரிவு}} (\emph{carivu}) is built from \ta{சரி}, \textquotedblleft{}to slip away, to slide down\textquotedblright{} (DEDR 2360), with the noun ending \emph{-vu}.
\end{cousinweb}

\begin{grammarlens}[title={What you've built}]
One more word for getting around.
\end{grammarlens}

\subsection*{Your turn}
\begin{itemize}
\raggedright
  \item \emph{Say it:} \emph{carivu}
  \item \emph{Say it:} \emph{carivu}, once more
  \item \emph{Say it:} \emph{malaic carivil āṭukaḷ mēyntukoṇṭiruntaṉa.} — Goats were grazing on the hillside.
  \item \emph{Recall:} say the Tamil for to pull over, then the Tamil for a curve, then say \emph{carivu} again
\end{itemize}

\begin{tcolorbox}[breakable,colback=teal!4,colframe=teal!35!black,title={Before you move on}]
What does \ta{சரிவு} mean? (\textquotedblleft{}A slope; a decline\textquotedblright{}.) Say it once more.
\end{tcolorbox}
\section[\texorpdfstring{\ta{ஏற்றம்} (ēṟṟam)}{ஏற்றம் (ēṟṟam)}]{\ta{ஏற்றம்} (ēṟṟam) — an ascent; a rise}

\label{lesson:TA-C408-erram}



\begin{quote}
Before the new one: say the Tamil for a curve, then the Tamil for a slope.
\end{quote}

\subsection*{You'll want to know: \ta{ஏற்றம்}}
\textbf{\ta{ஏற்றம்}} — \emph{ēṟṟam} — \textquotedblleft{}an ascent; a rise\textquotedblright{}.

\begin{quote}
\ta{இந்த} \ta{ஏற்றத்தில்} \ta{மிதிவண்டி} \ta{ஓட்டுவது} \ta{கடினம்}.
\end{quote}

\begin{quote}
\emph{inta ēṟṟattil mitivaṇṭi ōṭṭuvatu kaṭiṉam.}
\end{quote}

\begin{quote}
Riding a bicycle up this climb is hard.
\end{quote}

Its opposite is \ta{இறக்கம்}, a descent.

\begin{cousinweb}
\textbf{\ta{ஏற்றம்}} (\emph{ēṟṟam}) is built from \ta{ஏறு}, \textquotedblleft{}to climb, to rise\textquotedblright{} (DEDR 916), with the noun ending \emph{-am} and the doubled stop.
\end{cousinweb}

\begin{grammarlens}[title={What you've built}]
One more word for getting around.
\end{grammarlens}

\subsection*{Your turn}
\begin{itemize}
\raggedright
  \item \emph{Say it:} \emph{ēṟṟam}
  \item \emph{Say it:} \emph{ēṟṟam}, once more
  \item \emph{Say it:} \emph{inta ēṟṟattil mitivaṇṭi ōṭṭuvatu kaṭiṉam.} — Riding a bicycle up this climb is hard.
  \item \emph{Recall:} say the Tamil for a curve, then the Tamil for a slope, then say \emph{ēṟṟam} again
\end{itemize}

\begin{tcolorbox}[breakable,colback=teal!4,colframe=teal!35!black,title={Before you move on}]
What does \ta{ஏற்றம்} mean? (\textquotedblleft{}An ascent; a rise\textquotedblright{}.) Say it once more.
\end{tcolorbox}
\section[\texorpdfstring{\ta{கவிழ்} (kaviḻ)}{கவிழ் (kaviḻ)}]{\ta{கவிழ்} (kaviḻ) — to overturn; to capsize}

\label{lesson:TA-C408-kavil}



\begin{quote}
Before the new one: say the Tamil for a slope, then the Tamil for an ascent.
\end{quote}

\subsection*{You'll want to know: \ta{கவிழ்}}
\textbf{\ta{கவிழ்}} — \emph{kaviḻ} — \textquotedblleft{}to overturn; to capsize\textquotedblright{}.

\begin{quote}
\ta{வளைவில்} \ta{ஒரு} \ta{லாரி} \ta{கவிழ்ந்தது}.
\end{quote}

\begin{quote}
\emph{vaḷaivil oru lāri kaviḻntatu.}
\end{quote}

\begin{quote}
A lorry overturned on the bend.
\end{quote}

Past (it …): \ta{கவிழ்ந்தது}, \emph{kaviḻntatu}.

Of a boat, to capsize: \ta{படகு} \ta{கவிழ்ந்தது}.

\begin{cousinweb}
\textbf{\ta{கவிழ்}} (\emph{kaviḻ}) is native Dravidian (DEDR 1335), \textquotedblleft{}to be turned bottom upwards, to capsize\textquotedblright{}.
\end{cousinweb}

\begin{grammarlens}[title={What you've built}]
One more word for getting around.
\end{grammarlens}

\subsection*{Your turn}
\begin{itemize}
\raggedright
  \item \emph{Say it:} \emph{kaviḻ}
  \item \emph{Say it:} \emph{kaviḻ}, once more
  \item \emph{Say it:} \emph{vaḷaivil oru lāri kaviḻntatu.} — A lorry overturned on the bend.
  \item \emph{Recall:} say the Tamil for a slope, then the Tamil for an ascent, then say \emph{kaviḻ} again
\end{itemize}

\begin{tcolorbox}[breakable,colback=teal!4,colframe=teal!35!black,title={Before you move on}]
What does \ta{கவிழ்} mean? (\textquotedblleft{}To overturn; to capsize\textquotedblright{}.) Say it once more.
\end{tcolorbox}
\section[\texorpdfstring{\ta{சுழல்} (cuḻal)}{சுழல் (cuḻal)}]{\ta{சுழல்} (cuḻal) — to spin; to whirl}

\label{lesson:TA-C408-culal}



\begin{quote}
Before the new one: say the Tamil for an ascent, then the Tamil for to overturn.
\end{quote}

\subsection*{You'll want to know: \ta{சுழல்}}
\textbf{\ta{சுழல்}} — \emph{cuḻal} — \textquotedblleft{}to spin; to whirl\textquotedblright{}.

\begin{quote}
\ta{சேற்றில்} \ta{சக்கரம்} \ta{வேகமாகச்} \ta{சுழன்றது}.
\end{quote}

\begin{quote}
\emph{cēṟṟil cakkaram vēkamākac cuḻaṉṟatu.}
\end{quote}

\begin{quote}
The wheel spun fast in the mud.
\end{quote}

Past (it …): \ta{சுழன்றது}, \emph{cuḻaṉṟatu}.

\ta{ல்} meets the past -\ta{ந்த}- as \ta{ன்ற}: \ta{சுழன்றது}, it spun.

\begin{cousinweb}
\textbf{\ta{சுழல்}} (\emph{cuḻal}) is native Dravidian (DEDR 2698), \textquotedblleft{}to whirl, to spin, to turn on an axis\textquotedblright{}.
\end{cousinweb}

\begin{grammarlens}[title={What you've built}]
One more word for getting around.
\end{grammarlens}

\subsection*{Your turn}
\begin{itemize}
\raggedright
  \item \emph{Say it:} \emph{cuḻal}
  \item \emph{Say it:} \emph{cuḻal}, once more
  \item \emph{Say it:} \emph{cēṟṟil cakkaram vēkamākac cuḻaṉṟatu.} — The wheel spun fast in the mud.
  \item \emph{Recall:} say the Tamil for an ascent, then the Tamil for to overturn, then say \emph{cuḻal} again
\end{itemize}

\begin{tcolorbox}[breakable,colback=teal!4,colframe=teal!35!black,title={Before you move on}]
What does \ta{சுழல்} mean? (\textquotedblleft{}To spin; to whirl\textquotedblright{}.) Say it once more.
\end{tcolorbox}
